%------------------------------ Ejercicio 1B.1 -------------------------------

\setcounter{eje}{0}
\begin{ejer}
Defina $f \colon [0,1] \to \mathbb{R}$ como sigue:
$$
f(a) =%
\begin{cases} 
0, %
&\text{si $a \noten \mathbb{Q}$.} \\
\frac{1}{n}, %
&\text{si $a=\frac{m}{n}$ para algunos $m\en\mathbb{N}_0,n\en\mathbb{Z}^+$ primos relativos.}
\end{cases}
$$
Muestre que $f$ es Riemann integrable y compute $\int_0^1 f$.
\end{ejer}
\begin{sol}
Para cualquier partición $P$ y todo intervalo $[x_{i-1},x_i]$ de $P$, existe al menos un irracional en $[x_{i-1},x_i]$ (densidad de $\mathbb{R}\setminus\mathbb{Q}$ en $\mathbb{R}$), así que existe un $x \en [x_{i-1},x_i]$ tal que $f(x) = 0$. Como además $f \ge 0$, entonces necesariamente $\inf_{[x_{i-1},x_i]}f = 0$. Como esto último ocurre en cada intervalo de la partición arbitraria $P$, se sigue que $L(f,P,[0,1]) = 0$ para toda partición $P$ de $[0,1]$; en consecuencia, $L(f,[0,1]) = 0$.\par
Por otro lado, para cada $n\en\mathbb{N}$, definimos $P_n$ como la \emph{partición regular} con longitud de intervalo constante igual a $\nicefrac{1}{n^3}$, es decir, $P_n$ está definida por los puntos $\set[\big]{\nicefrac{k}{n^3} \given 0 \leq k \leq n^3}$. Ahora bien, fijemos un $n\en\mathbb{N}$ arbitrario. Notemos que hay una cantidad finita de puntos en los que $f$ es mayor o igual a $\nicefrac{1}{n}$, ya que si $a\en[0,1]$ cumple $f(a) \geq \nicefrac{1}{n}$, entonces forzosamente $a\en\mathbb{Q}$ y $a=\nicefrac{m}{k}$ para algunos $m\en\mathbb{N}_0,k\en\mathbb{Z}^+$ primos relativos; de esto se sigue que $\nicefrac{1}{k} = f(a) \geq \nicefrac{1}{n}$, así que $k \leq n$; como además $a\en[0,1]$, tenemos que $m \leq k$. Esto prueba que, si $f(a) \geq \nicefrac{1}{n}$, entonces $a = \nicefrac{m}{k}$ con $1 \leq k \leq n$ y $0 \leq m \leq k$; en otras palabras, 
\[
f^{-1}\big(\big[\tfrac{1}{n},\infty\big)\big)%
\subseteq%
\set*{\tfrac{m}{k} \given 1 \leq k \leq n, 0 \leq m \leq k }
\]
Este último conjunto tiene, a lo sumo, tantos elementos como tiene el conjunto que lo indexa, como el conjunto de índices está contenido en $\set{1,\dots,n} \times\set{0,\dots,n}$, se tiene
\[
\operatorname{\#}\big(
    f^{-1}\big(\big[\tfrac{1}{n},\infty\big)\big)
\big)%
\;\leq\;%
\operatorname{\#}\big(
    \set*{\tfrac{m}{k} \given 1 \leq k \leq n, 0 \leq m \leq k }
\big)%
\;\leq\;%
n(n+1)
\]

Ahora acotemos la suma superior: la partición $P_n$ se puede escribir como unión disjunta de dos conjuntos, $P^1_n$ y $P^2_n$, de modo que $P^1_n$ consta de aquellos intervalos de $P_n$ que tienen algún punto de $f^{-1}\big([\nicefrac{1}{n},\infty)\big)$ y $P^2_n$ de aquellos que no. Como cada punto de $f^{-1}\big([\nicefrac{1}{n},\infty)\big)$ puede estar en, a lo sumo, dos intervalos distintos, se sigue que $\operatorname{\#}(P^1_n) \leq 2n(n+1)$. Luego, gracias a que la longitud de todos los intervalos es $\nicefrac{1}{n^3}$:
\[
U(f,P_n,[0,1]) = \sum_{I\en P^1_n}\frac{1}{n^3}\sup\nolimits_If +\sum_{I\en P^2_n}\frac{1}{n^3}\sup\nolimits_If
\]
Pero, por definición, $f \leq \nicefrac{1}{n}$ sobre todo intervalo $I\en P^2_n$, es decir, $\sup_I f \leq \nicefrac{1}{n}$ para cada $I\en P^2_n$. De modo similar, $f\leq 1$ en todo su dominio, así que siempre tenemos $\sup_I f \leq 1$ para cualquier intervalo. Luego
\[
U(f,P_n,[0,1])%
= \sum_{I\en P^1_n}\frac{1}{n^3} +\sum_{I\en P^2_n}\frac{1}{n^3}\frac{1}{n}%
= \operatorname{\#}(P^1_n)\frac{1}{n^3} + \operatorname{\#}(P^2_n)\frac{1}{n^4}
\]
pero ya vimos que $\operatorname{\#}(P^1_n) \leq 2n(n+1)$ y también $\operatorname{\#}(P^2_n) \leq \operatorname{\#}(P_n) = n^3$. Por lo tanto:
\[
U(f,P_n,[0,1])%
\leq 2\frac{n(n+1)}{n^3} + \frac{n^3}{n^4}%
= \frac{3}{n} + \frac{2}{n^2}%
\]
Por lo tanto se tiene, para cada $n\en\mathbb{N}$, que 
\[
0 = L(f,[0,1]) \leq U(f,[0,1]) \leq U(f,P_n,[0,1]) \leq \frac{3}{n} + \frac{2}{n^2}
\]
como $\lim_{n\to\infty}\frac{3}{n} + \frac{2}{n^2} = 0$, el teorema de compresión de límites implica que 
\[
0 = L(f,[0,1]) = U(f,[0,1])
\]
así que $f$ es Riemann integrable y $\int^1_0f = 0$.
\end{sol}

%------------------------------ Ejercicio 1B.2 -------------------------------

\begin{ejer}
Suponga que $f \colon [a, b] \to \mathbb{R}$ es acotada. Pruebe que $f$ es Riemann integrable si y solo si
$$
L(-f, [a, b]) = -L(f, [a, b]).
$$
\end{ejer}
\begin{sol}
Si $P$ es cualquier partición de $[a,b]$ dada por $x_0<x_1<\cdots <x_n$, entonces
\[
L(-f,P,[a,b])=\sum_{i=1}^n(x_i-x_{i-1})\inf_{{[x_{i-1},x_i]}} (-f),
\]
pero, para cada sub-intervalo $I_i=[x_{i-1},x_i]$, se cumple $\inf_{I_i}(-f)=-\sup_{I_i}f$. Por tanto,
\[
L(-f,P,[a,b])%
=\sum_{i=1}^n(x_i-x_{i-1})\inf_{{[x_{i-1},x_i]}} (-f)
=-\sum_{i=1}^n(x_i-x_{i-1})\sup_{{[x_{i-1},x_i]}} f
= -U(f,P,[a,b])
\]
Luego, si denotamos por $\mathcal{P}$ al conjunto de particiones de $[a,b]$, podemos tomar supremos:
\[
\sup_{P\en\mathcal{P}} L(-f,P,[a,b])
=\sup_{P\en\mathcal{P}} \bigl(-U(f,P,[a,b])\bigr)
=-\inf_{P\en\mathcal{P}} U(f,P,[a,b])
\]
Es decir, esto se traduce, por definición, en $L(-f,[a,b]) = -U(f,[a,b])$. Esta igualdad implica la siguiente cadena de equivalencias:
\begin{align*}
f\text{ es integrable }%
\quad&\iff\quad
U(f,[a,b])=L(f,[a,b])\\
\quad&\iff\quad
-U(f,[a,b])=-L(f,[a,b])\\
\quad&\iff\quad
L(-f,[a,b])=-L(f,[a,b])%
\end{align*}
Esto concluye la prueba.
\end{sol}

%------------------------------ Ejercicio 1B.4 -------------------------------

\setcounter{eje}{3}
\begin{ejer}
Dé un ejemplo de funciones acotadas $f, g\colon [0,1] \to \mathbb{R}$ tales que
$$
L(f, [0,1]) + L(g, [0,1]) < L(f + g, [0,1])
$$
y
$$
U(f + g, [0,1]) < U(f, [0,1]) + U(g, [0,1]).
$$
\end{ejer}
\begin{sol}
Tomemos $f = \chi_{\mathbb{Q} \cap[0,1]}$ la función característica de $\mathbb{Q} \cap[0,1]$ y $g = \chi_{[0,1]\setminus\mathbb{Q}}$ la función característica de $[0,1]\setminus\mathbb{Q}$. La imagen de ambas funciones es exactamente $\set{0,1}$, así que ambas son acotadas y no negativas.\par
Ahora bien, para cualquier sub-intervalo $I\subset [0,1]$, la densidad de $\mathbb{Q}$ y de $\mathbb{R\setminus\mathbb{Q}}$ implican que tanto $I\cap\mathbb{Q}$ como 
$I\cap([0,1]\setminus\mathbb{Q})$ son no vacíos. Por tanto,
\[
\inf\nolimits_I f=0,\qquad \sup\nolimits_I f=1,\qquad 
\inf\nolimits_I g=0,\qquad \sup\nolimits_I g=1.
\]
De aquí se sigue que, para toda partición $P$ de $[0,1]$,
\[
L(f,P,[0,1])=0,\qquad U(f,P,[0,1])=1,
\]
y análogamente
\[
L(g,P,[0,1])=0,\qquad U(g,P,[0,1])=1.
\]
Tomando supremos e ínfimos sobre todas las particiones,
\[
L(f,[0,1])=0,\qquad U(f,[0,1])=1,
\qquad
L(g,[0,1])=0,\qquad U(g,[0,1])=1.
\]
Por otro lado, para todo $x\in[0,1]$ se tiene
\[
(f+g)(x)=\chi_{\mathbb{Q}}(x)+\chi_{[0,1]\setminus\mathbb{Q}}(x)=1.
\]
Es decir, $f+g\equiv 1$ en $[0,1]$. En consecuencia, para cualquier partición $P$,
\[
L(f+g,P,[0,1])=\sum_{i=1}^n 1\cdot (x_i-x_{i-1})=1,
\qquad
U(f+g,P,[0,1])=1,
\]
y por tanto
\[
L(f+g,[0,1])=U(f+g,[0,1])=1.
\]
Finalmente,
\[
L(f,[0,1])+L(g,[0,1])=0+0=0<1=L(f+g,[0,1]),
\]
y
\[
U(f+g,[0,1])=1<2=U(f,[0,1])+U(g,[0,1]).
\]
Esto prueba lo pedido.
\end{sol}